% commentaire.tex — f170r, vv. 635–678 (conversion initiale des travaux antérieurs)
\comm{637–40}{κἄπειτ᾽ ἀποδύσετ᾽ ἐννέα παίδων μητέρα; ... καὶ νὴ Δία τίτθος γ᾽ ὥσπερ ἡμεῖς οὐκ ἔχει.}{Le manuscrit de Ravenne porte en marge du vers 638 μν\auDessus{η}{σ}, indiquant ainsi que Mnésiloque prend la parole ; le manuscrit ne présente alors plus aucun signe de changement de locuteur jusqu'au vers 641\nc{Il est d'usage dans notre manuscrit que le \italique{signe} soit un παράγραφος semblable à un tiret cadratin ` --- ' lorsque le changement de locuteur se fait en début de vers. Si le changement intervient en milieu de vers, il y a plus souvent un \latin{punctum duplex} ` : ' ou bien une abréviation indiquant le nom du personnage qui prend la parole. Par abus de langage, nous utiliserons le terme de \italique{paragraphos} pour désigner indifféremment l'un de ces signes.}. Or Mnésiloque ne peut aucunement prononcer le vers 639, commentaire sur sa poitrine étonnamment plate (ὡς καὶ στιβαρά τις φαίνεται καὶ καρτερά).
Il y a donc un \italique{paragraphos} à rétablir.
Je dis \italique{un} car les éditeurs suivent quasi-unanimement Brunck en `rétablissent' deux : au début du vers 638 (Clistène) et au début du vers 639 (la femme -- Mika -- reprend la parole).
Une autre solution me semble plus légère : laisser à Mnésiloque les deux vers 637-638 et à Mika les vers 639-640.
Cette retouche comporte surtout l'avantage d'être bien meilleure quant à l'action et au jeu de scène.\\
Presque tous les éditeurs depuis Brunck, disais-je, retirent le vers 638 à Mnésiloque pour l'accorder à Clisthène. Double erreur ! \ \ \ \italique{Primo}, parce qu'ils sont alors contraints de surajouter une seconde émendation : transformer le χάλᾳ du manuscrit en un χάλα (impératif lancé par Clisthène) alors que le iota adscrit est nettement écrit. \ \ \ \italique{Secundo}, parce que Mika, au vers 636, avait déjà ordonné à Clisthène de déshabiller Mnésiloque (ἀπόδυσον αὐτόν) ; il n'y a aucune de raison de douter qu'il ne soit en train de lui obéir. Un peu d'action, que diable ! une pièce d'Aristophane ne doit pas avoir la même pesanteur que l'administration française. (Selon nos éditeurs, Mika ordonnerait à Clisthène de déshabiller Mnésiloque qui transmettrait alors l'ordre à l'intéressé... qui finirait par ne rien faire ?..)\\
Le vers 638, χαλᾷ ταχέως τὸ στρόφιον ὦναίσχυντε σύ, était pourtant excellent dans la bouche de Mnésiloque : Clisthène est en train de le déshabiller, il feint l'indignation (χαλᾷ est un subjonctif dit \italique{délibératif}) devant l'empressement (ταχέως) que cet homme -- car Clisthène reste un homme -- met à enlever le \italique{soutien gorge} d'une femme. « Il défait mon strophion avec promptitude. Impudent ! » Mnésiloque continue de jouer une femme -- vertueuse même, mais un peu tard. Qu'y a-t-il de plus limpide ? l'apostrophe ὦναίσχυντε σύ conclue logiquement la tirade de mère effarouchée qu'affecte le peu crédible Mnésiloque (et qui fait d'ailleurs écho au précédent ἀναίσχυντός τις εἶ (v. 611) lancé par le même Mnésiloque au même Clisthène\nc{Quelques éditeurs avaient même prétexté de ce double ἀναίσχυντος pour attribuer l'ἀναίσχυντός τις εἶ du vers 611 aussi à Clisthène, histoire que leurs multiples tripatouillages présentent une apparence d'homogénéité...}).}
\comm{641}{στερίφη γάρ εἰμι.}{Jeu de mots d'Aristophane : comment Mnésiloque explique-t-il la platitude de sa poitrine ? tout simplement parce qu'il a les chairs \italique{fermes} (sens premier de στέριφος, le Scythe l'emploie avec cette acception au vers 1185). Mais στέριφος peut aussi avoir le sens dérivé de \italique{stérile} ; sens que le spectateur ne privilégiera qu'en entendant la seconde partie de la réplique de Mnésiloque (κοὐκ ἐκύησα). C'est la raison pour laquelle Aristophane recourt à un mot rare, στέριφος. (Le scholiaste de noter : ἀντὶ τοῦ στείρα, commentaire qui se retrouve dans de nombreux lexiques, cf. l'apparat des témoins\nc{Et pour l'apparat le plus fourni, se reporter à celui de l'édition Theodoridis de Photius.}.)}
\comm{642–3}{νῦν· τότε δὲ μήτηρ ἦσθα παίδων ἐννέα. / ἀνίστασ᾽ ὀρθός. ποῖ τὸ πέος ὠθεῖς κάτω;}{Il n'y a aucun \italique{paragraphos} aux vers 642 et 643.
Pourtant, ce n'est assurément pas Mnésiloque (qui a la parole au vers précédent) qui peut prononcer ces vers. On ajoute souvent deux changement d'interlocuteurs : Mika pour le vers 642 et Clisthène pour le vers 643. Par simplicité, je préfère n'en ajouter qu'un et attribue donc le distique entier à Clisthène.}
\comm{643}{ἀνίστασ᾽ ὀρθός.}{«~Redresse-toi !~» L'injonction adressée à Mnésiloque est bien-sûr à double entente, et Clisthène se fera doublement obéir. Mais interrogeons-nous : pourquoi Mnésiloque s'était-il assis ou recroquevillé ? réponse simple : pour cacher sa virilité. Ils se redressent. Mnésiloque tente une dernière fois de cacher ce qui doit être caché et Clisthène : ποῖ τὸ πέος ὠθεῖς κάτω; «~Où repousses-tu ce \italique{péos} ?!~» Le \italique{péos} passe alors au verso de Mnésiloque, d'où peut le contempler Mika, admiratrice. Avant que ne s'ensuive la scène burlesque où Clisthène paraît fort intéressé aussi.\\
Mnésiloque s'était habillé dans cet ordre : κροκωτός, στρόφιον et ἔγκυκλον. Le στρόφιον et l'ἔγκυκλον viennent d'être retirés, mais pas le κροκωτός (au vers 941, Mnésiloque le porte toujours) ; comment se fait-il que son \italique{olisbos} soit apparent malgré la présence de cette robe safranée ? parce que la crocote est composée d'un tissu très fin qui la rend transparente\nc{Cf. le cratère à figures rouges (ca. -380) représentant une scène de comédie : un vieillard entouré de deux gardes et vêtu de la crocote... qui laisse tout transparaître ; il est reproduit en annexe. Cratère «~St. Agata Antigone~», St. Agata dei Goti I, collection Rainone-Mustilli. Inventorié par Arthur D. Trendall , \italique{Phlyax Vases}, «~BICS~» suppl. LXIX (1967).}}
\comm{643}{Note sur τὸ ἔγκυκλον $=$ \italique{une jupe}.}{Clisthène vient de retirer le στρόφιον de Mnésiloque, ce qui entraîna quelques remarques sur sa poitrine. Or, lors de la scène d'habillement (vers 253-261), Mnésiloque avait successivement enfilé le κροκωτός, le στρόφιον \italique{puis, à la fin,} l'ἔγκυκλον. Il résulte donc que l'ἔγκυκλον \italique{est un vêtement qui ne peut pas recouvrir la poitrine}. (Sinon Clisthène aurait dû commencer par l'encycle avant le στρόφιον.) 
Ce ne peut donc être en aucun cas un \italien{mantello ampio, pesante}\nc{Comme le conclue Franca Perusino dans son article «~\italien{L'ἔγκυκλον, un mantello femminile nelle commedie di Aristofane}~», \italique{Quaderni Urbinati di Cultura Classica}, 72.3 (2002), pp. 129-133. Ses conclusions erronées ont malheureusement été reproduites dans la plupart des travaux portant sur le costume de l'Ancienne comédie... espérons que cette note rectificatrice mettra enfin un terme à la contagion !}.
Allongeons quelque peu cette note afin que soit irréfragable l'identification de l'ἔγκυκλον avec une sorte de jupe courte.\\
Que pourrait-il être d'autre ce vêtement qui sitôt retiré laisse tout émoustillé le mari de Praxagora ; voyez plutôt (\italique{Assemblée des femmes}, vers 535-538) :\citComm{Βλέπυρος -- εἶτ' οὐ τὸ σαυτῆς ἱμάτιον ἐχρῆν σ' ἔχειν;\\
ἀλλ' ἔμ' ἀποδύσασ' ἐπιβαλοῦσα \textbf{τοὔγκυκλον}\\
ᾤχου καταλιποῦσ' ὡσπερεὶ προκείμενον,\\
μόνον οὐ στεφανώσασ' οὐδ' ἐπιθεῖσα λήκυθον.\\
Alors que n'as-tu pris ton propre manteau [pour sortir] ? Mais tu me dévêts et jettes sur moi \italique{ton encycle}, puis tu t'en vas, me laissant là comme un cadavre exposé, hors que tu as omis de me couronner et de placer près de moi un lécythe.}Le champ lexical est celui d'une exposition funèbre\nc{Cf. par exemple \italique{Lysistrata}, vv. 599 sqq.} comme la traduction de Willems le laisse entendre, mais le λήκυθον du vers 538 doit être pris \italique{sensu obsceno} selon la matrice métaphorique : λήκυθος $=$ ἀλάβαστος\nc{Cf. la saillie du chaud-lapin Kinésias soumis à la torture par Myrrhine (\italique{Lysistrata}, v. 947) :\citComm{\textbf{Μυρρίνη} -- λαβὲ τόνδε τὸν ἀλάβαστον. --- \textbf{Κινησίας} -- ἀλλ᾽ ἕτερον ἔχω.}} $=$ \latin{phallus}.
Blépyros reproche à sa femme de ne pas s'être acquitté du \italique{nome de mâtines} (ὄρθριος νόμος\nc{\italique{Assemblée des femmes}, vers 741. Par jeu de mots avec l'ὄρθιος νόμος, le \italique{mode orthien}, qu'Hésychius (ο 1188) n'a pas manqué de nous définir :\citComm{ὄρθιον νόμον· οὕτως ὠνομάσθαι φασὶν αὐτὸν ἀπὸ τοῦ ῥυθμοῦ, καθάπερ καὶ τὸν τροχαῖον.}(Cf. aussi sur le mode orthien \italique{Cavaliers} 1280 ; \italique{Acharniens} 16.)\\
Les allusions d'Aristophane à cette tradition matinale à laquelle sacrifient ses personnages sont légions ; ne citons que \italique{Lysistrata} v. 966 puisque le passage est encore sous nos yeux.}).
Ce vêtement permet à Aristophane un jeu de mot avec l'\italique{enceinte} de la ville fortifiée de Pylos\nc{\italique{Lysistrata}, vers 1162-1163. Pour conclure la paix, l'ambassadeur laconien demande à récupérer l'\italique{encyle} :\citComm{Λάκων -- ἁμές γε λῶμες, αἴ τις ἁμὶν τὤγκυκλον\\
λῇ τοῦτ᾽ ἀποδόμεν. --- Λυσιστράτη -- ποῖον ὦ τᾶν; --- Λάκων -- τὰν Πύλον,\\
ἇσπερ πάλαι δεόμεθα καὶ βλιμάττομες.}} -- oui ! --, la Πύλος, cette ville qui déclenche intarissablement la verve éroticomique de notre poète, jonglant avec l'autre matrice métaphorique (féminine cette fois) : πύλος $=$ πύλη $=$ θύρα $=$ \latin{pudenda muliebria} (raison pour laquelle l'Athénien ne veut pas la céder). 
Il y avait certainement à l'époque de la représentation de \italique{Lysistrata} -- sans doute d'une année antérieure à celle des \italique{Thesmophories} -- des tractations pour que la paix fût conclue en échange de la remise de Pylos ou, par synecdoque, de son \italique{encycle}. 
Ce n'est pas autrement qu'il faut comprendre le passage au début de la même pièce (vers 111-114) : Lysistrata demande : \citComm{ἐθέλοιτ᾽ ἂν οὖν, εἰ μηχανὴν εὕροιμ᾽ ἐγώ,\\
μετ᾽ ἐμοῦ καταλῦσαι τὸν πόλεμον;\\
Consentirez-vous donc, si j'en trouve le moyen, à mettre fin avec moi à la guerre ?}et Aristophane de répondre, par la voix d'une femme :\citComm{\decalage{μετ᾽ ἐμοῦ καταλῦσαι τὸν πόλεμον;}νὴ τὼ θεώ\\
ἔγωγ᾽ ἂν οὖν κἂν εἴ με χρείη \textbf{τοὔγκυκλον}\\
τουτὶ καταθεῖσαν ἐκπιεῖν αὐθημερόν.\\
Oui, par les deux déesses. J'en suis, quand il me faudrait déposer l'\italique{encycle} que voici en gage pour boire dès aujourd'hui [à la paix : les traités, σπονδαί, sont en grec, littéralement, des \italique{libations}].}Il faut donc bien-sûr aussi attribuer à cette même femme -- et comme ne le font pas les éditeurs d'Aristophane -- le vers 104 :\citComm{ὁ δ᾽ ἐμός γε τελέους ἑπτὰ μῆνας ἐν Πύλῳ.}(Elle se dit prête à faire tomber l'\italique{encycle} dès le retour de son mari en garnison depuis sept longs mois à Pylos.)
Aristophane, par deux fois dans la même pièce, exhorte ses contemporains à renoncer à Pylos, c'est dire l'importance qu'il accordait à l'éventualité de la paix ! On connaît la suite...
Ce vêtement, donc, qui enclot le sadinet, puisqu'il faut parler net... Alors ? que peut-il bien être ce vêtement ? S'il vous plaît, ne me répondez pas de chœur avec tous les fonctionnaires de la recherche : «~un lourd et pesant manteau~».\\
Mais délaissons un peu nos primes amours, les comédies d'Aristophane, pour aborder un terrain plus aride et cependant qui n'est certainement pas à négliger, je veux parler de la lexicologie.
Ce sens de \italique{jupe} pour l'ἔγκυκλον est confirmé, ce me semble, par le commentaire le plus précis : celui du scholiaste de \italique{Lysistrata} (vers 113), qui glose ἔγκυκλον ainsi :\citComm{τὸ περιβόλαιον τὸ ζωνάριον.}J'interprète ζωνάριον comme un adjectif se rapportant à περιβόλαιον\nc{Chantraine, page 74 de sa \italique{Formation des noms en grec ancien}, précise : «~la finale -άριον qui a fourni essentiellement des diminutifs s'observe parfois dans des adjectifs, cf. σωμα παιδαριον dans des inscription de Delphes~».} : \italique{le vêtement / la couverture}, celle qui enceint, \italique{au niveau de la taille} ; une jupe donc. Citons aussi ce passage du copieux commentaire à l'\italique{Iliade} (976.13) d'Eustathe de Thessalonique (il commente \italique{Iliade} XIV, v. 178) :\citComm{ἰστέον δὲ ὅτι, καθάπερ ὁ ἑανὸς οὕτω καὶ τὸ ἔκκυκλον ἢ ἔγκυκλον, οὗ καὶ ὁ κωμικὸς μέμνηται, γυναικεῖον ἦν ποτε φόρημα, φησὶ γοῦν Παυσανίας ὅτι ἔγκυκλον περιπόρφυρον\nc{On lit en épigraphie : ἔνκυκλον περιποίκιλο, ἔγκυκλον ποικίλον. À l'inverse, on croise aussi un ἔνκυκλον λευκ[όν].} ἱμάτιον καὶ χιτὼν γυναικεῖος, ὃν ἔνδοθεν ἐνδύονται γυναῖκες, εἶτα τὸ ἔνδυμα. λέγεται δὲ φησίν [sic, je lis : σφῇσιν] ὑπ' ἐνίων ἐπιγονατίς, διότι μέχρι γονάτων διήκει.\\
Il faut savoir que, de même que l'ἑανὸς [= riche robe de femme dans le passage d'Homère qu'Eustathe commente], il y a aussi l'ἔκκυκλον (ou ἔγκυκλον), dont s'est aussi souvenu le Comique [= Aristophane]. C'était un vêtement (φόρημα) féminin porté jadis. Ce qui est sûr, c'est que Pausanias indique que l'ἔγκυκλον était un vêtement (ἱμάτιον) aux franges pourpres (περιπόρφυρον) et une tunique (χιτὼν) féminine, que les femmes revêtaient à l'intérieur, et ensuite un habit (ἔνδυμα). Il est qualifié par elles-mêmes, d'après certains, d'ἐπιγονατίς, parce qu'il atteignait les genoux.}Ce vêtement érotique d'intérieur peut-il être autre chose qu'une \italique{jupe sexy} ?}
\comm{644}{τάλαν.}{Ce mot, le vocatif de τάλας, tant masculin que féminin, n'est mis par Aristophane que dans la bouche de femmes\nc{\italique{Lysistrata}, vers 102, 910 et 914 ; \italique{Grenouilles}, vers 559 ; \italique{Assemblée des femmes}, vers 124, 658 et 1005 ; \italique{Ploutos}, vers 706 et 1055.}.}
\comm{647–8}{ἰσθμόν τιν᾽ ἔχεις ὦνθρωπ᾽· ἄνω τε καὶ κάτω / τὸ πέος διέλκεις πυκνότερον Κορινθίων.}{Clisthène fait référence au \grec{δίολκος} : un chemin dallé de 8 kilomètres qui permettait de faire traverser aux bateaux l'isthme de Corinthe grâce à une sorte de rail rudimentaire et d'éviter ainsi le long et périlleux contournement du Péloponnèse\nc{Voir à ce sujet l'article de M. J. T. Lewis, « Railways in the Greek and Roman world », in \italique{Early Railways : a Selection of Papers from the First International Early Railways Conference}, 2001, p. 8-19. Thucydide (III, 15 ; VIII, 7) évoque le δίολκος dans sa \italique{Guerre du Péloponnèse}.}. Pendant la scène burlesque de déshabillage qui vient d'avoir lieu, Mnésiloque, pour cacher sa virilité – et certainement aussi émoustillé d'être effeuillé de la sorte –, a actionné à plusieurs reprises le mécanisme permettant de lever et baisser son \italique{phallus}. Admiratif devant de telles performances, Clisthène le flatte d'être encore mieux outillé que les machines qui permettaient le transbordement des bateaux.\\
Voilà pour l'image `τὸ πέος διέλκεις πυκνότερον Κορινθίων' ! il reste à expliquer la première partie de la réplique de Clisthène : ἰσθμόν τιν᾽ ἔχεις, «~tu possèdes un [drôle de] genre d'\italique{isthme}~».
ἰσθμός désigne proprement un \italique{passage étroit} et, par analogie, le \italique{cou}. Le spectateur athénien saisissait instamment l'image : par \italique{isthme} il fallait entendre \italique{pénis} grâce à la matrice métaphorique suivante : ἰσθμός = αὐχήν\nc{αὐχήν a le sens de πέος par exemple au vers 681 de \italique{Lysistrata} : le chœur des vieillards a maille à partir avec l'engeance féminine ; après avoir rappelé par une métaphore hippique combien la femme est portée à la luxure, il en vient à la conclusion que le meilleur moyen pour dominer ces fougueuses juments est de les chevaucher -- la fameuse \italique{phallocratie} ? Il formule cette idée à l'aide d'une image carcérale : il faut assujettir leur \italique{bois percé} (= \italique{carcan}) au \italique{cou} (αὐχήν) :
\citComm{ἀλλὰ τούτων χρῆν ἁπασῶν ἐς τετρημένον ξύλον\\
ἐγκαθαρμόσαι λαβόντας \textbf{τουτονὶ τὸν αὐχένα}.\\
Eh bien, il importe de s'assurer d'elles toutes, et d'assujettir dans un carcan \italique{ce cou-là}.}
On peut avoir un autre exemple d'αὐχήν avec ce sens-là au vers 184 du \italique{Cyclope} d'Euripide.} = τράχηλος\nc{τράχηλος est phallossémique au vers 20 des \italique{Grenouilles}. Aristophane commence la pièce avec un rapide survol des blagues lourdes et vulgaires de ses concurrents aux concours comiques ; ce n'est pas sa muse qui s'abaisserait à de si grossières manœuvres ! non, \italique{on n'en soufflera mot}, οὐκ ἐρεῖ ! οὐκ ἐρεῖ !.. et puis au vers 19 :
\citComm{ὢ τρισκακοδαίμων ἄρ᾽ \textbf{ὁ τράχηλος οὑτοσί},\\
ὅτι θλίβεται μέν, τὸ δὲ γέλοιον οὐκ ἐρεῖ.\\
Ô trois fois malencontreuse \italique{nuque}, qui est moulue, et n'a pas le mot pour rire !}
Il faut bien-sûr lire -- je dis «~bien-sûr~», mais personne ne l'avait proposé jusqu'ici... --, il faut donc lire :
\citComm{ἄρ᾽ ὁ τράχηλος οὑτοσί,\\
ὅτι θλίβεται μέν, τὸ δὲ γέλοιον \textbf{οὐ κε ῥεῖ;}\\
Est-ce que le τράχηλος que voici, puisqu'il est pressé, pour la farce ne devrait pas couler ?}} = πέος.}
\comm{651}{εἰσεκύλισα.}{Homère emploie le verbe κυλίνδειν en parlant de malheurs que le ciel \italique{fait rouler} sur quelqu'un : πῆμα θεὸς Δαναοῖσι κυλίνδει (\italique{Il}. 17, 688)\nc{Le moyen κυλίνδεσθαι se dit de malheur lui-même «~roulant~» sur quelqu'un : νῶιν δὴ τόδε πῆμα κυλίνδεσθαι, \italique{Il.}, 11, 347 ; cf. \italique{Od}. 2, 163 ; 8, 81.}. La même image se trouve chez Aristophane : νὴ τὸν Διόνυσον ἄπορά γ᾽ ἡμῖν πράγματα / δαίμων τις ἐσκεκύκληκεν ἐς τὴν οἰκίαν (\italique{Guêp.} 1474 s.), « ah! les inextricables tracas qu'un dieu a fait dévaler sur notre maison. » Le même verbe est employé d'une manière un peu différente, \italique{Thesm}. 767, mais toujours à propos de malheurs : ὁ μὲν γὰρ αἴτιος / κἄμ᾽ ἐσκυλίσας εἰς τοιαυτὶ πράγματα / οὐ φαίνεταί πω, « le coupable, celui qui m'a fait rouler dans de tels malheurs, n'apparaît toujours pas. » Cf. \italique{Thesm}. 650 s. : κακοδαίμων ἐγώ, / εἰς οἷ᾽ ἐμαυτὸν εἰσεκύλισα πράγματα. Il est clair que, dans ce cas, ἐμβάλλειν τινὰ εἰς συμφοράς n'est qu'un synonyme expressif\nc{Connu d'autres auteurs ; Phérécr. 146 [K = 156 K.-A.] : εἰκῆ μ᾿ ἐπῆρας ὄντα τηλικουτονὶ / πολλοῖς ἐμαυτὸν ἐγκυλῖσαι πράγμασιν ; Cf. Xén. \italique{Mém.} I, 2, 22 : τοὺς εἰς ἔρωτας ἐγκυλισθέντας ; Ἀπολλόδωρος ὁ Καρύστιος 5, 7 [ $=$ Athénée 280e] : εἰκῆ τέ πως / ἡμᾶς κυλίνδους᾽ ὅντιν᾽ ἂν τύχῃ τρόπον (sc. τύχη ἄγροικος).} de ἐμβάλλειν (p. ex. ἐμβάλλειν τινὰ εἰς συμφοράς, Antiphon 3, 4, 10 ; εἰς ἀτυχίας, Eschin. 3, 79 ; εἰς κακόν, Ar. \italique{Nuées} 1460). (Note de Taillardat, §558.)}
\comm{655–86}{Le chœur compris entre les vv. 655-686.}{La \nieme{19} note d'Alphonse Willems est d'une telle richesse d'explication que je ne peux que la reproduire intégralement ci-dessous. Pour ce qui touche à la colométrie, je reprends celle de Lætitia\nc{\italique{The songs of Aristophanes}, pp. 418-422.}. (Le passage des vers 667-686 a de fortes correspondances métriques avec les vers 707-725).\\
Ce petit chœur n'a pas été examiné avec assez d'attention. On n'a pas remarqué qu'il s'accompagne d'un jeu de scène unique dans l'œuvre d'Aristophane. Trois morceaux le composent : \\
\ \ \ 1° Les vv. 655 à 658 appartiennent au coryphée. Mnésiloque n'est peut-être pas le seul homme qui ait pénétré parmi les femmes. Il importe de s'en assurer, et, après avoir déposé les manteaux et allumé les torches, de « tourner et virer tout à l'entour du Pnyx, et d'inspecter avec soin les tentes et les passages ».\\
\ \ \ 2° Les quatre suivants (659-662) sont des trochées, et le changement de mètre indique un changement d'interlocuteur. C'est, en effet, la réponse du \italique{parastate} ou chef du demi-chœur [le \italique{Ravennas} porte sur sa marge, η\auDessus{μ}{χ}ι γ\auDessus{υ}{ν}, écrite d'une seconde main, mais dont l'indication se retrouve sur le manuscrit de Münich]. En disant qu'il y avait lieu d'inspecter le Pnyx, les tentes et les passages, la coryphée indiquait clairement qu'il fallait quitter l'orchestre, c'est-à-dire l'enceinte sacrée adjacente au temple\nc{On peut conclure de là que le Thesmophorion, qu'on savait situé sur une hauteur, l'était sur la colline du Pnyx et voisin du lieu d'assemblée.}. La parastate annonce donc que le chœur va « s'élancer dehors (ἐξορμᾶν) d'un pied léger\nc{Le simple, ὀρμᾶν, qui se lit plus loin, v. 953, et \italique{Gren}., 478, signifierait « prendre sa course ».} ». Puis elle engage celle qu'elle appelle τὴν πρώτην à commencer tout de suite la reconnaissance. Comment faut-il interpréter τὴν πρώτην ? D'y voir, comme on le fait généralement, une locution adverbiale (\latin{quam primum} ou \latin{prima quaque via}) me semble inadmissible. Avec M. Mazon, je crois que ce terme désigne une personne déterminée\nc{\italique{Essai sur la composition des comédies d'Aristophane}, Paris (1904), p. 131, n. 4.} ; suivant moi, ce serait τὴν πρώτην ἱέρειαν, appelée aussi τὴν προστάτιν, savoir la première, il y en avait beaucoup\nc{Il y en avait deux par dème (Isée, \italique{Succession de Kiron}, 19-20), donc plus de deux cents.}, des prêtresses temporaires, ἀρήτειραι, de Déméter, celle qui, tout à l'heure, avait prononcé les prières d'usage (vv. 295 et suiv.), en un mot, la coryphée. 
Callimaque, parlant d'une de ces prêtresses, ἱερέη Δήμητρος, se sert d'une désignation analogue : πολλῶν προστασίη νέων γυναικῶν\nc{\italique{Anthol. Palat}., VII, 728.}. \\
\ \ \ 3° Le morceau qui suit (vv. 663-684) est un \italique{hyporchême}. Une partie du chœur reste en place et chante, tandis que l'autre, en quête de quelque complice de Mnésiloque\nc{εἴ τις ἐν τόποις ἑδραῖος. Qu'avons-nous affaire d'ajouter, avec MM. Blaydes et van Leeuwen, τούτοισι τοῖς devant τόποις ? La même locution adverbiale se lit dans un fragment de Nausicratès : [τρίγκλαι] ἃς κλύδων Αἰξωνικὸς ἐν τόποις παιδεύεται (\italique{Fragm. Com. Gr}., t. IV, p. 575 [= Athénée 325e = Nausicratès, fr. 1.6-11 Kassel-Austin ; le texte n'est pas établi avec certitude : on trouve aussi : ἐντόπους à la place de ἐν τόποις. Même si l'exemple n'est pas probant, Willems n'en a pas moins raison ; Lætitia remarque que l'adjonction de Blaydes accroît la disparité entre les deux stances des vers 663-666 -- là où l'on chercherait à rétablir une symétrie.]) et la locution inverse se lit dans Sophocle : ἥτις μ᾽ ἀποικίσειεν ἐκ τόπων (\italique{Trach}., 955), où jamais personne n'a trouvé à redire (ἐκ τούτων τῶν τόπων.). Je ne vois pas que ces tournures soient plus bizarres que ἐν ποσί et ἐκ ποδῶν, ou bien que ἐν χώρᾳ et ἐκ χώρας, πρὸ ou παρὰ χειρῶν, etc.}, évolue silencieusement, σιωπῇ (v. 660), va et vient, entre et sort, cherchant partout, dans l'orchestre et dehors, καὶ τὰ τῇδε καὶ τὰ δεῦρο (v. 666), jusqu'à ce qu'il déclare n'avoir rien trouvé (687-688).\\}
\comm{656}{τῶν θ᾽ ἱματίων ἀποδύσας.}{Willems, fidèle à son poste de \italique{gardien de la tradition}, justifie contre les retoucheurs de texte l'expression ἀποδῦναι $+$ gén. (C'est sa \nieme{20} note aux \italique{Thesmophories}.)}
\comm{657}{εἰσελήλυθεν.}{La leçon du manuscrit de Ravenne est \italique{contra metrum}. Tannegui Le Fèvre (Faber dans l'apparat), le père de Mme Dacier, avait proposé la simple correction ἐσελήλυθε. Austin\nc{Voir ses notes : « Textual problems in Ar. \italique{Thesm}. » (1987) ; je reprends la trame de celle qui nous intéresse présentement.} la réfuta : ἐς ne précède jamais une voyelle chez Aristophane. Restent les conjectures de Fritzsche et de Handley. Fritzsche proposa de lire ἀνελήλυθε, l'expliquant par une haplographie ἀνὴρ (ἀν)ἐλήλυθε suivie d'une complétion postérieure εἰσ-. Handley proposa ἐπελήλυθε (la confusion entre ΙC et Π en onciales étant une erreur courante) -- hypothèse beaucoup plus simple et, partant, plus satisfaisante.}
